% Directly authored; companion to gamma_optimization_full_note.md.
% Build from the repository root with latexmk -pdf -outdir=/tmp/gamma-full-latex.
\documentclass[11pt]{article}
\usepackage[T1]{fontenc}
\usepackage{amsmath,amssymb,amsthm}
\usepackage{newtxtext,newtxmath}
\usepackage[letterpaper,margin=0.78in]{geometry}
\usepackage{graphicx,microtype,xcolor,caption,fancyhdr}
\usepackage[colorlinks=true,linkcolor=blue!40!black,urlcolor=blue!50!black]{hyperref}
\definecolor{ink}{HTML}{263442}
\DeclareMathOperator{\sech}{sech}
\DeclareMathOperator{\csch}{csch}
\DeclareMathOperator{\sign}{sign}
\DeclareMathOperator{\diag}{diag}
\theoremstyle{definition}
\newtheorem{proposition}{Proposition}
\newtheorem{lemma}{Lemma}
\newtheorem{theorem}{Theorem}
\newtheorem{corollary}{Corollary}
\newcommand{\repo}{https://github.com/samlayton99/precision-mlps/blob/experiment/frozen-gamma-probe}
\newcommand{\evidence}{\repo/results/checkpoint_D_optimizers/expD36_frozen_gamma_probe/full_sweep/refinements/gamma_factorized_kernel}
\graphicspath{{results/checkpoint_D_optimizers/expD36_frozen_gamma_probe/full_sweep/refinements/gamma_factorized_kernel/}}
\hypersetup{pdftitle={From gamma to slow target directions: the complete argument},
 pdfsubject={Explicit gamma attenuation, target alignment, full proofs, and verified learning times}}
\setlength{\parindent}{0pt}
\setlength{\parskip}{5pt plus 1pt}
\setlength{\emergencystretch}{2em}
\setlength{\headheight}{14pt}
\captionsetup{font=small,labelfont=bf,justification=raggedright,singlelinecheck=false}
\pagestyle{fancy}
\fancyhf{}
\fancyhead[L]{\small\color{ink}Gamma, target alignment, and readout learning}
\fancyhead[R]{\small Full technical note}
\fancyfoot[C]{\small\thepage}
\renewcommand{\headrulewidth}{0.3pt}

\begin{document}
\thispagestyle{plain}
{\color{ink}\fontsize{23}{26}\selectfont\bfseries From gamma to slow target directions:\\the complete argument\par}
\vspace{8pt}
Small gamma smooths the features. To turn this into a statement about
optimization, we must show that some part of the target lies in directions on
which the resulting kernel acts weakly. This note supplies that missing step,
proves the learning-time consequences, and explains how retaining the full
kernel yields the sharp predictions in the experiments.

There are two levels of conclusion. An explicit bound on high-frequency kernel
action guarantees slow target energy and a learning delay under stated target
and geometry conditions. That bound can be conservative. A calculation
retaining the whole gamma-filtered kernel gives much sharper, target-specific
intervals. The measured factor of 986.59 belongs to the second calculation;
it is not the sharpness of the first bound.

\textbf{Notation.} Throughout this note, $\mu_i$ denotes a kernel eigenvalue.
The symbol $\lambda=\gamma h$ is reserved for dimensionless bandwidth when the
centers have spacing $h$. The physical slope is $\gamma$. The step size is
$\eta_\gamma$, the largest kernel eigenvalue is $L_\gamma$, and
$a_i=\eta_\gamma\mu_i$ is a per-update decay rate. Norms are Euclidean for
vectors and spectral for matrices. Every additional symbol is defined where
it enters. Earlier notes used $\lambda_i$ for eigenvalues and $\beta$ for
bandwidth; this note supersedes that notation.

\section{The model and the optimization question}
Fix training inputs $x_1,\ldots,x_m$, centers $c_1,\ldots,c_W$, and a nonzero
target. Only the raw readout coefficients are trained:
\begin{equation*}
 f_{\gamma,\theta}(x)=b+\sum_{j=1}^Ww_j\tanh(\gamma(x-c_j)),\qquad
 \theta=(b,w_1,\ldots,w_W)^T.
\end{equation*}
Normalize both the samples of the features and the target by $\sqrt m$:
\begin{equation*}
 (J_\gamma)_{i0}=m^{-1/2},\quad
 (J_\gamma)_{ij}=m^{-1/2}\tanh(\gamma(x_i-c_j)),\quad
 y_i=m^{-1/2}f^\star(x_i).
\end{equation*}
The loss is $\frac12\|J_\gamma\theta-y\|^2$, exactly half empirical mean
squared error. Define
\begin{equation*}
 K_\gamma=J_\gamma J_\gamma^T,\qquad L_\gamma=\|K_\gamma\|,\qquad\theta_0=0.
\end{equation*}
The bias column has norm one, so $L_\gamma\ge1$. Use a prescribed stable step
$0<\eta_\gamma L_\gamma\le1$. Our archived experiments have
$\eta_\gamma L_\gamma\simeq1/2$. Changing gamma holds the target, inputs,
centers, raw coefficient metric, and step-size rule fixed.

Representability asks whether some coefficients fit the target. Optimization
asks how many updates starting at zero are needed to find sufficiently accurate
coefficients. A target can be representable while the relevant kernel
eigenvalues are very small.

\subsection*{An exact example with no spectral ambiguity}
Take one center at zero, samples $-s,+s$ with $s>0$, and target values $-1,+1$.
Put $d=\tanh(\gamma s)$. Then
\begin{equation*}
 K_\gamma=\frac12\begin{pmatrix}1+d^2&1-d^2\\1-d^2&1+d^2\end{pmatrix}.
\end{equation*}
The mean direction $(1,1)/\sqrt2$ has eigenvalue $1$. The contrast direction
$(-1,1)/\sqrt2$, which is exactly the target direction, has eigenvalue $d^2$.
These statements follow by multiplying the matrix by the two vectors. Thus
$L_\gamma=1$, while the target's normalized eigenvalue is
$\tanh^2(\gamma s)$.

Every positive gamma fits exactly with $b=0,w=1/d$. Nevertheless, at step
$\eta_\gamma=\kappa\in(0,1)$,
\begin{equation*}
 E_\gamma(n)=(1-\kappa\tanh^2(\gamma s))^n,\qquad
 n_\epsilon=\left\lceil\frac{\log(1/\epsilon)}
 {-\log(1-\kappa\tanh^2(\gamma s))}\right\rceil.
\end{equation*}
As $\gamma s\to0$, $n_\epsilon\sim\log(1/\epsilon)/(\kappa\gamma^2s^2)$.
A cap $\gamma\le\bar\gamma$ gives the exact lower bound obtained by replacing
gamma with $\bar\gamma$ in the displayed time; it is attained at the cap.
With $\kappa=1/2$ and $\epsilon=0.01$, the counts are 925 for $\gamma s=0.1$
and 14 for $\gamma s=1$. A constant target instead decays as $(1-\kappa)^n$,
independently of gamma. This proves both the possible optimization obstruction
and the necessity of a target condition.

\section{What determines training: rates weighted by the target}
Let $r_n=y-J_\gamma\theta_n$ and $E_\gamma(n)=\|r_n\|/\|y\|$.
The acquisition time is
\begin{equation*}
 n_\epsilon(\gamma)=\inf\{n\in\mathbb N_0:E_\gamma(n)\le\epsilon\},
\end{equation*}
with value $+\infty$ if the set is empty. A 1\% relative residual means a
relative squared residual of $10^{-4}$.

\begin{proposition}[Exact GD evolution]
Let $u_i$ be an orthonormal eigenbasis of $K_\gamma$, including its nullspace,
with eigenvalues $\mu_i\ge0$. Set
\begin{equation*}
 a_i=\eta_\gamma\mu_i,\qquad p_i=|u_i^Ty|^2/\|y\|^2.
\end{equation*}
Then $\sum_i p_i=1$, $0\le a_i\le1$, and
\begin{equation}\label{eq:gd}
 \boxed{r_n=(I-\eta_\gamma K_\gamma)^ny,\qquad
 E_\gamma(n)^2=\sum_i p_i(1-a_i)^{2n}.}
\end{equation}
\end{proposition}
\begin{proof}
The gradient is $J_\gamma^T(J_\gamma\theta-y)$. Hence
$\theta_{n+1}=\theta_n+\eta_\gamma J_\gamma^Tr_n$ and
$r_{n+1}=(I-\eta_\gamma K_\gamma)r_n$. Since $r_0=y$, induction proves
the first identity. Expand $y$ in the orthonormal eigenbasis and square its
norm to obtain the second. The convention at $n=0$ is that every mode
contributes its initial energy, including a mode with $a_i=1$. Each factor
is nonincreasing, so the residual is nonincreasing.
\end{proof}
The limiting squared residual is $E_\infty^2=\sum_{\mu_i=0}p_i$, which is at
most every finite-time squared residual. This is the projection of the target
outside the feature span; it does not compare the tanh model with a
step-function model.

For $\eta_\gamma=\kappa/L_\gamma$, the rates are
$a_i=\kappa\mu_i/L_\gamma$. Overall kernel rescaling cancels. The relevant
object is the distribution of target energy across normalized eigenvalues,
rather than the condition number alone. Individual eigenvectors need not be
tracked across gamma; within repeated eigenspaces their choice is arbitrary,
but total target energy in each spectral subspace is well defined.

Define the \textbf{slow target mass} at rate cutoff $t\in(0,1)$ by
\begin{equation*}S_\gamma(t)=\sum_{a_i\le t}p_i.\end{equation*}
\begin{corollary}[Slow mass implies a delay]
\begin{equation}\label{eq:massbound}
 E_\gamma(n)\ge\sqrt{S_\gamma(t)}(1-t)^n.
\end{equation}
If $\sqrt{S_\gamma(t)}>\epsilon$, then
\begin{equation}\label{eq:masstime}
 n_\epsilon(\gamma)\ge\left\lceil
 \frac{\log(\sqrt{S_\gamma(t)}/\epsilon)}{-\log(1-t)}\right\rceil.
\end{equation}
\end{corollary}
\begin{proof}
Keep only the nonnegative terms with $a_i\le t$ in \eqref{eq:gd}, and use
$(1-a_i)^{2n}\ge(1-t)^{2n}$. Taking square roots gives
\eqref{eq:massbound}; solving its necessary inequality for a crossing gives
\eqref{eq:masstime}.
\end{proof}
This corollary is already target aware. We now connect its slow mass to gamma
without assuming Fourier waves are finite-kernel eigenvectors.

\section{Gamma is an exact smoothing operation on the features}
For a bounded function $\varphi:\mathbb R\to\mathbb R$, define the unit-mass
density and smoothing operation
\begin{equation*}
 \rho_\gamma(t)=\frac\gamma2\sech^2(\gamma t),\qquad
 (\mathcal S_\gamma\varphi)(x)=\int_{\mathbb R}\rho_\gamma(x-t)\varphi(t)\,dt.
\end{equation*}
Here $\varphi$ may be a constant or a step feature.
The calligraphic $\mathcal S_\gamma$ distinguishes this operation from the
scalar slow mass $S_\gamma(t)$.
\begin{lemma}[Exact smoothing and multiplier]
Constants are preserved, and
$\mathcal S_\gamma\sign(\cdot-c)=\tanh(\gamma(\cdot-c))$.
For $k_{\rm step}(t,s)=1+\sum_j\sign(t-c_j)\sign(s-c_j)$,
\begin{equation}\label{eq:convolution}
 \begin{split}
 k_\gamma(x,x')&=\iint\rho_\gamma(x-t)\rho_\gamma(x'-s)k_{\rm step}(t,s)\,dt\,ds,\\
 (K_\gamma)_{i\ell}&=k_\gamma(x_i,x_\ell)/m.
 \end{split}
\end{equation}
With convention $\widehat\rho(\omega)=\int\rho(t)e^{-i\omega t}dt$, the multiplier is
\begin{equation}\label{eq:multiplier}
 \boxed{M_\gamma(\omega)=\frac{z}{\sinh z},\qquad
 z=\frac{\pi|\omega|}{2\gamma},\qquad M_\gamma(0)=1.}
\end{equation}
\end{lemma}
\begin{proof}
The density has cumulative integral $[1+\tanh(\gamma t)]/2$. Splitting the
convolution at $c$ gives
$2\int_{-\infty}^{x-c}\rho_\gamma(u)du-1=\tanh(\gamma(x-c))$.
Unit mass preserves the bias. Insert the finite sum defining $k_{\rm step}$
to prove \eqref{eq:convolution}. Bounded features and an integrable density
justify the exchanges; the step functions need not belong to whole-line $L^2$.

For the transform, put $\xi=\omega/\gamma>0$ and integrate
$e^{-i\xi z}/\cosh^2z$ around a rectangle with imaginary heights zero and
$\pi$. The vertical integrals vanish as its real endpoints tend to infinity.
The upper horizontal integral is $-e^{\pi\xi}$ times the lower one. The
double pole at $z_0=i\pi/2$ has residue $i\xi e^{\pi\xi/2}$: near $z_0$,
$1/\cosh^2z=-1/(z-z_0)^2+O(1)$. Consequently,
\begin{equation*}
 (1-e^{\pi\xi})\int_{\mathbb R}\frac{e^{-i\xi u}}{\cosh^2u}\,du
 =-2\pi\xi e^{\pi\xi/2},
\end{equation*}
so the integral is $\pi\xi/\sinh(\pi\xi/2)$. Scaling $u=\gamma t$ and
dividing by two gives \eqref{eq:multiplier}. Evenness handles negative
frequencies and continuity gives the value at zero.
\end{proof}
For $z>0$, $(z/\sinh z)'=(\sinh z-z\cosh z)/\sinh^2z<0$: the numerator
starts at zero and has derivative $-z\sinh z$. Thus smaller gamma gives
stronger attenuation at a fixed nonzero frequency. A cap
$\gamma\le\bar\gamma$ implies $M_\gamma(\omega)\le M_{\bar\gamma}(\omega)$,
and $M_\gamma(\omega)\sim2ze^{-z}$ at high frequency.

Equation \eqref{eq:convolution} smooths the first and second input arguments
separately. It is an identity before sampling, not multiplication of the
finite sample matrix by an unspecified discrete convolution matrix. The
identity alone does not order the finite kernels or their target-weighted spectra.

\section{Keep the finite geometry: an explicit matrix construction}
We need a finite representation in which the gamma-dependent factors are
known and the discarded part is bounded. This construction approximates the
original nonperiodic tanh features; it does not make the training problem periodic.

Let $R_0=\max_{i,j}|x_i-c_j|$, choose a half-period $T>R_0$, and choose a
positive integer $Q$. For $\ell=1,\ldots,Q$, define
\begin{equation*}
 \omega_\ell=(2\ell-1)\pi/T,\quad b_\ell=4/[\pi(2\ell-1)],\quad
 h_{\gamma,Q}(z)=\sum_{\ell=1}^Q b_\ell M_\gamma(\omega_\ell)\sin(\omega_\ell z).
\end{equation*}
\begin{lemma}[Feature remainder]
Put $d=\pi^2/(2\gamma T)$ and $u=d(2Q+1)$. On $|z|\le R_0$,
\begin{equation}\label{eq:featureerror}
 \begin{split}
 |\tanh(\gamma z)-h_{\gamma,Q}(z)|&\le e_{\gamma,Q},\\
 e_{\gamma,Q}&:=4e^{-2\gamma(T-R_0)}+
 \frac{4\pi}{\gamma T}\frac{e^{-u}}{(1-e^{-2u})(1-e^{-2d})}.
 \end{split}
\end{equation}
\end{lemma}
\begin{proof}
The periodic square wave $\sign(\sin(\pi z/T))$ agrees with $\sign(z)$ on
$(-T,T)$. Their difference has magnitude at most two outside. After
smoothing, its magnitude is at most
$1-\tanh(\gamma(T-z))+1-\tanh(\gamma(T+z))$, bounded by the first term of
\eqref{eq:featureerror}. The square wave has odd sine coefficients $b_\ell$;
smoothing multiplies each by \eqref{eq:multiplier}, giving an absolutely
convergent series. For odd $k$, the smoothed coefficient equals
$2\pi\csch(dk)/(\gamma T)$. At $k\ge2Q+1$, use
$\csch(dk)\le2e^{-dk}/(1-e^{-2u})$ and sum the geometric series with ratio
$e^{-2d}$. This proves the second term.
\end{proof}
Define three matrices, giving their entries so their roles are explicit:
\begin{itemize}
\item $F$ has $m$ rows and $2Q+1$ columns:
$1,\sin(\omega_1x),\cos(\omega_1x),\ldots$, sampled at the inputs and divided
by $\sqrt m$. It contains only the sample geometry.
\item $C$ has $2Q+1$ rows and $W+1$ columns. Its bias row is
$(1,0,\ldots,0)$, and its remaining bias-column entries are zero. For hidden
column $j$, $C_{2\ell-1,j}=b_\ell\cos(\omega_\ell c_j)$ and
$C_{2\ell,j}=-b_\ell\sin(\omega_\ell c_j)$. It contains only the center
geometry and step coefficients.
\item $D_\gamma$ is diagonal with entries
$1,M_\gamma(\omega_1),M_\gamma(\omega_1),\ldots$.
It contains all the gamma dependence.
\end{itemize}
\begin{proposition}[Finite factorization and kernel error]
\begin{equation}\label{eq:factor}
 \widetilde J_\gamma=FD_\gamma C,\qquad
 \widetilde K_\gamma=FD_\gamma CC^TD_\gamma F^T.
\end{equation}
With $\delta_\gamma=\sqrt W e_{\gamma,Q}$,
\begin{equation}\label{eq:kernelerror}
 \begin{split}
 \|J_\gamma-\widetilde J_\gamma\|&\le\delta_\gamma,\\
 \|K_\gamma-\widetilde K_\gamma\|&\le
 \Delta_\gamma:=(2\|\widetilde J_\gamma\|+\delta_\gamma)\delta_\gamma.
 \end{split}
\end{equation}
\end{proposition}
\begin{proof}
Expand $\sin(\omega(x-c))=\sin(\omega x)\cos(\omega c)-\cos(\omega x)\sin(\omega c)$
to obtain \eqref{eq:factor}. The bias is exact, and each normalized hidden
column has error norm at most $e_{\gamma,Q}$. Summing squared column errors
bounds the Frobenius norm, hence the spectral norm, by $\delta_\gamma$.
Set $H=J_\gamma-\widetilde J_\gamma$ and expand
$K_\gamma-\widetilde K_\gamma=H\widetilde J_\gamma^T+\widetilde J_\gamma H^T+HH^T$
to obtain \eqref{eq:kernelerror}.
\end{proof}
The sampled waves in $F$ need not be orthogonal, and $CC^T$ need not be
diagonal. Therefore $M_\gamma(\omega)^2$ is not generally an individual
finite-kernel eigenvalue. Equation \eqref{eq:factor} retains both effects.
It constructs the prediction using fixed geometry and explicit multipliers;
it is not a fit to an optimization trajectory.

\section{The missing implication: attenuation forces slow target energy}
For any unit sample vector $v$,
\begin{equation}\label{eq:rayleigh}
 v^TK_\gamma v=\|J_\gamma^Tv\|^2=\sum_i\mu_i|u_i^Tv|^2.
\end{equation}
Thus small correlations with all features imply small spectral strength. The
following theorem turns that observation into a target-dependent statement,
including the effect of eigendirection rotation.

Choose a frequency cutoff $0<\Omega\le\omega_Q$. Let $P_{<\Omega}$ be the
orthogonal projector, in the empirical Euclidean norm, onto the span of the
bias and columns of $F$ with frequency below $\Omega$. Define
$R_\Omega=I-P_{<\Omega}$ and
\begin{equation*}\alpha=\|R_\Omega y\|/\|y\|.\end{equation*}
This is an explicit assumption on the target's component outside the sampled
low-frequency span. It does not define the target using unknown slow
eigenvectors. It also does not assert that a continuously high-frequency
function is exactly orthogonal on a finite grid.

Let $F_{\ge\Omega}$ collect the remaining wave columns and $C_{\ge\Omega}$
the matching rows. Put
\begin{equation*}
 B_\Omega=\|F_{\ge\Omega}\|^2\|C_{\ge\Omega}\|^2,\qquad
 A_\gamma=\eta_\gamma[B_\Omega M_\gamma(\Omega)^2+\Delta_\gamma].
\end{equation*}
Both the subspace and $B_\Omega$ are determined by fixed geometry,
$T,Q,\Omega$, independently of gamma and GD.

\begin{theorem}[Gamma-dependent slow-target guarantee]
For every $t\in(0,1)$,
\begin{equation}\label{eq:targetbridge}
 \boxed{\begin{aligned}
 S_\gamma(t)&\ge(\alpha-\sqrt{A_\gamma/t})_+^2,\\
 E_\gamma(n)&\ge(\alpha-\sqrt{A_\gamma/t})_+(1-t)^n,
 \end{aligned}}
\end{equation}
where $(z)_+=\max(z,0)$. If $b_\gamma(t):=(\alpha-\sqrt{A_\gamma/t})_+>\epsilon$,
then
\begin{equation}\label{eq:targettime}
 \boxed{n_\epsilon(\gamma)\ge\left\lceil
 \frac{\log(b_\gamma(t)/\epsilon)}{-\log(1-t)}\right\rceil.}
\end{equation}
\end{theorem}
\begin{proof}
\textbf{Step 1: gamma bounds action on the specified subspace.}
For a unit $v$ in the range of $R_\Omega$, all low-frequency coordinates of
$F^Tv$ vanish. Therefore
\begin{equation*}
 \begin{aligned}
 v^T\widetilde K_\gamma v
 &=\|C_{\ge\Omega}^TD_{\gamma,\ge\Omega}F_{\ge\Omega}^Tv\|^2\\
 &\le B_\Omega M_\gamma(\Omega)^2.
 \end{aligned}
\end{equation*}
The diagonal multiplier decreases with absolute frequency. Applying
\eqref{eq:kernelerror} gives $v^T(\eta_\gamma K_\gamma)v\le A_\gamma$.

\textbf{Step 2: weak action forces the direction into slow modes.}
Let $P_{\rm fast}$ project onto rates $a_i>t$ and
$P_{\rm slow}=I-P_{\rm fast}$. The spectral theorem gives
$\eta_\gamma K_\gamma\succeq tP_{\rm fast}$. Hence
$\|P_{\rm fast}v\|\le\sqrt{A_\gamma/t}$.

\textbf{Step 3: transfer that fact to the target.}
If $\alpha=0$, \eqref{eq:targetbridge} is trivial. Otherwise put
$\widehat y=y/\|y\|$ and $v=R_\Omega\widehat y/\alpha$.
Then $v^T\widehat y=\alpha$. Orthogonality of the two spectral projectors and
Cauchy--Schwarz give
\begin{equation*}
 \begin{aligned}
 \alpha&\le |(P_{\rm slow}v)^TP_{\rm slow}\widehat y|
       +|(P_{\rm fast}v)^TP_{\rm fast}\widehat y|\\
       &\le\sqrt{S_\gamma(t)}+\sqrt{A_\gamma/t}.
 \end{aligned}
\end{equation*}
Rearranging proves the slow-mass bound. Corollary 1 proves the residual and
time bounds. No shared eigenvectors across gamma, commuting operators, or
identification of waves with eigenvectors was used.
\end{proof}
\textbf{Interpretation.} Gamma enters the guarantee through the squared
multiplier. The target enters through $\alpha$, the step through
$\eta_\gamma$, and the finite geometry through $B_\Omega$ and the approximation
error. A positive statement requires $\alpha>\sqrt{A_\gamma/t}+\epsilon$.
The mere presence of a tiny high-frequency component does not force a delay
to a tolerance larger than that component.

\subsection*{A uniform common-slope cap corollary}
For an explicit uniform statement, fix $0<\gamma_-\le\gamma\le\bar\gamma$
and use the same geometry, $T,Q,\Omega$, and rule
$\eta_\gamma=\kappa/L_\gamma$ with $0<\kappa\le1$.
Let $V_{\bar\gamma,Q}$ be the second term of \eqref{eq:featureerror}
evaluated at $\bar\gamma$, and define
\begin{equation*}
 \begin{split}
 e_*&=4e^{-2\gamma_-(T-R_0)}+V_{\bar\gamma,Q},\qquad\delta_*=\sqrt W e_*,\\
 \Delta_*&=(2\sqrt{W+1}+\delta_*)\delta_*.
 \end{split}
\end{equation*}
Then Theorem 1 holds throughout this slope interval with
\begin{equation*}
 A_*:=\kappa[B_\Omega M_{\bar\gamma}(\Omega)^2+\Delta_*]
\end{equation*}
in place of $A_\gamma$. For any class of targets with $\alpha\ge\alpha_0$,
replace $\alpha$ by $\alpha_0$ to obtain a common learning-delay guarantee.
\begin{proof}
The distant-transition error is largest at $\gamma_-$. Every omitted Fourier
coefficient is $b_\ell M_\gamma(\omega_\ell)$ with $b_\ell>0$, so its
absolute sum is at most the sum at $\bar\gamma$, bounded by
$V_{\bar\gamma,Q}$. Thus the feature error is at most $\delta_*$.
Since $\|J_\gamma\|\le\|J_\gamma\|_F\le\sqrt{W+1}$, expanding the kernel
error around $J_\gamma$ gives the displayed $\Delta_*$. Finally
$L_\gamma\ge1$ implies $\eta_\gamma\le\kappa$, and multiplier monotonicity
gives $A_\gamma\le A_*$ when the kernel-error bound is taken as $\Delta_*$.
Apply Theorem 1.
\end{proof}
The positive lower slope endpoint makes this particular periodic approximation
uniform. This corollary does not assert a sharp bound over all
$0<\gamma\le\bar\gamma$, heterogeneous slopes, or varying center sets.
Its constants can be large. We have not established that this new coarse
corollary is tight on the reported experiment.

\section{Why the short argument is not automatically a sharp time law}
Theorem 1 loses information by bounding a product with separate operator
norms, replacing all high-frequency multipliers by one worst-case value, and
converting a single action bound into spectral mass. Corollary 1 then treats
every mode below the cutoff as if it decayed at that cutoff. These are valid
inequalities; none is generally an equality. Increasing $Q$ controls the
approximation error but does not remove these other sources of slack.

There is also no general matrix ordering from increasing the filter entries
alone. For an algebraic example, take
\begin{equation*}
 G=\begin{pmatrix}1&1\\1&1\end{pmatrix},\qquad
 D=\diag(1,d),\quad 0<d<1.
\end{equation*}
Although the diagonal entries of $D$ are at most those of the identity,
$\det(G-DGD)=-(1-d)^2<0$. Thus $G-DGD$ is indefinite. For
$v=(1,-1)/\sqrt2$, $v^TGv=0$ but $v^TDGDv=(1-d)^2/2$.
This is a counterexample to an inference from diagonal filtering alone, not
a claim about the particular tanh geometry in the experiment. Target
alignment and geometric coupling cannot be discarded.

Under additional simultaneous diagonalization, the simpler statement does
hold: if an orthonormal frequency basis diagonalizes both the reference
kernel and the smoothing operator, then the filtered kernel eigenvalues are
$M_\gamma(\omega)^2\mu_{\rm step}(\omega)$ in that basis. Multiplying the
two diagonal operators proves this immediately. Our finite model is not
assumed to have this property. Equally spaced centers alone do not remove
boundaries, sampling effects, or aliasing.

\section{Sharp timing by retaining the entire filtered kernel}
We now keep all couplings and all target weights in \eqref{eq:factor}. Let
$\widetilde K_\gamma\widetilde u_i=\widetilde\mu_i\widetilde u_i$ denote
its positive modes, $\xi_i=\widetilde u_i^Ty$, and
$y_\perp=y-\sum_i\xi_i\widetilde u_i$. Set
$q_i=1-\eta_\gamma\widetilde\mu_i$ and assume both the true and approximate
kernels satisfy the stable-step condition at the prescribed step. Define
\begin{equation}\label{eq:approxcurve}
 \widetilde E_\gamma(n)^2=
 \frac{\|y_\perp\|^2+\sum_i|\xi_i|^2q_i^{2n}}{\|y\|^2}.
\end{equation}
This is GD in the approximate filtered feature model. It is not a projection
of the original GD residual onto a different basis. Proposition 1 proves its formula.

\begin{theorem}[Transfer of the full curve]
Write $Z=K_\gamma-\widetilde K_\gamma$ and let $\|Z\|\le\Delta_\gamma$.
Then $|E_\gamma(n)-\widetilde E_\gamma(n)|\le d_n$, where
\begin{equation}\label{eq:transfer}
 d_n=\min\left\{n\eta_\gamma\Delta_\gamma,
 \frac{\eta_\gamma}{\|y\|}
 \left[n\|Zy_\perp\|+
 \sum_i|\xi_i|\|Z\widetilde u_i\|\sum_{j=0}^{n-1}q_i^j\right]\right\}.
\end{equation}
\end{theorem}
\begin{proof}
Put $U=I-\eta_\gamma K_\gamma$ and $V=I-\eta_\gamma\widetilde K_\gamma$.
Both have norm at most one. The identity
\begin{equation*}
 U^n-V^n=\sum_{j=0}^{n-1}U^{n-1-j}(U-V)V^j
\end{equation*}
is obtained by expanding the sum and canceling adjacent terms. Since
$U-V=-\eta_\gamma Z$, contraction bounds the norm of the sum by
$\|U^n-V^n\|\le n\eta_\gamma\Delta_\gamma$.
Alternatively, apply the identity to $y$ and substitute
$V^jy=y_\perp+\sum_i\xi_iq_i^j\widetilde u_i$. The triangle inequality
and contraction of $U$ give the second expression in \eqref{eq:transfer}.
Finally, $|\|U^ny\|-\|V^ny\||\le\|(U^n-V^n)y\|$. Divide by $\|y\|$ and
take the smaller bound.
\end{proof}
The geometric sum equals $(1-q_i^n)/(1-q_i)$ for a positive rate, and its
continuous limit is $n$ at zero rate. The first bound uses the analytic
feature remainder; the second also uses the original kernel's action on
known model directions. Neither needs optimizer iterates. Actions can be
evaluated as
$Zv=(J-\widetilde J)(\widetilde J^Tv)+J((J-\widetilde J)^Tv)$ without
subtracting nearly equal full kernel products.

\begin{corollary}[Necessary and sufficient times]
Put $\ell(n)=\max(0,\widetilde E_\gamma(n)-d_n)$ and
$u(n)=\min(1,\widetilde E_\gamma(n)+d_n)$.
If integers $n_-<n_+$ satisfy $\ell(n_-)>\epsilon$ and $u(n_+)\le\epsilon$,
then
\begin{equation}\label{eq:interval}
 n_-+1\le n_\epsilon(\gamma)\le n_+.
\end{equation}
\end{corollary}
\begin{proof}
The lower test excludes the threshold at $n_-$ and, by monotonicity, every
earlier step. The upper test exhibits a step attaining it. The upper envelope
need not itself be monotone; only the checked witness is used.
\end{proof}
Bounds valid at different $Q$ may be intersected. For positive finite
intervals $L_g\le n_\epsilon(g)\le U_g$ and $L_h\le n_\epsilon(h)\le U_h$,
division gives $L_g/U_h\le n_\epsilon(g)/n_\epsilon(h)\le U_g/L_h$.

This route preserves the complete target-weighted spectrum. It explains the
accurate timing calculation without suggesting that the coarse guarantee in
Theorem 1 has the same tightness.

\section{The intermediate spectral evidence and the executed training}
The primary experiment fixes 559 hidden features, 8,193 equally spaced training
inputs on $[-1,1]$, and equispaced centers with a halo. The target is
\begin{equation*}
 f^\star(x)=\sin(2\pi x)+\tfrac12\sin(6\pi x)+\tfrac14\sin(10\pi x).
\end{equation*}
All readouts start at zero. Gamma is 8, 12, 16, or 64, and the saved GD steps
satisfy $\eta_\gamma L_\gamma\simeq0.5$. The Fourier construction uses $T=8$
and up to $Q=2048$ odd harmonics. Its target weights come from a rectangular
SVD of $FD_\gamma C$, with no fitted training rates.

At rate cutoff $t=10^{-6}$, the numerically computed slow target masses are
4.33333\%, 0.00107533\%, 0.0000285067\%, and 0.000148493\%, respectively.
These are fractions of squared target norm expressed as percentages, not
residual percentages. The cutoff was chosen as an explanatory diagnostic
after the experiment, not as a precommitted selection criterion.

\begin{figure}[htbp]
 \centering
 \includegraphics[width=\linewidth]{spectral_bridge.pdf}
 \caption{\textbf{The missing link is target energy in slow modes.}
 A shows $S_\gamma(t)$ from the gamma-filtered finite kernels; dots mark
 $t=10^{-6}$. The horizontal guide is target energy $10^{-4}$, or 0.01\%,
 corresponding to a 1\% residual. B compares the full-spectrum residual
 predictions with the lower bounds
 $\sqrt{S_\gamma(10^{-6})}(1-10^{-6})^n$; its horizontal guide is a 1\%
 residual and its vertical guide is 16,013 updates. All curves here are
 calculations from the archived kernels, not executed trajectories. The
 spectra and cutoff bounds are FP64 diagnostics, not independently
 interval-certified spectral masses.}
\end{figure}

For gamma 8, Corollary 1 evaluated with the computed mass gives
$E_8(16013)\gtrsim0.20486$; the full-spectrum prediction is $0.22967$.
At gamma 64 the predicted residual is $0.00999980$. The same cutoff gives a
gamma-8 necessary time of 3,035,752 updates, versus the full prediction and
executed crossing of 15,798,313. This visible looseness is exactly why the
final timing calculation retains more spectral information. A single cutoff
does not order all intermediate gammas: at this cutoff gamma 16 has less mass
than gamma 64, despite taking longer to reach 1\%.

\begin{figure}[htbp]
 \centering
 \includegraphics[width=\linewidth]{pi_brief_three_panel.pdf}
 \caption{\textbf{Connect the mechanism to executed optimization.}
 A evaluates the analytic multiplier; frequency guides mark the three target
 frequencies. B shows kernel-predicted GD curves as lines, actual saved GD
 residuals as open circles, and measured 1\% crossings as diamonds. The
 dotted verticals locate predicted crossings. C connects actual Adam
 checkpoints under the same optimizer settings across gammas. The GD proof
 does not apply to Adam. Both optimizers start at zero. GD measurements are
 subsampled by update count for visibility, with every saved residual retained
 in the evidence.}
\end{figure}

The executed GD crossing counts are 15,798,313; 186,057; 61,792; and 16,013.
The corresponding independently certified combined intervals are
[15,784,048, 15,812,623], [186,057, 186,058], [61,792, 61,792], and
[16,013, 16,013]. Their gamma-8/gamma-64 ratio is bounded by
985.70--987.49, versus a measured 986.59. All 1,796 saved GD residuals agree
with the $Q=2048$ filtered-model predictions within $8.92\times10^{-15}$ in
absolute relative residual. This numerical agreement is distinct from
interval certification.

The analytic-remainder-only gamma-8 interval is [15,783,830, 15,812,843];
sharpness therefore does not require the action refinement. A control
preserving the gamma-64 eigenvectors and relative eigenvalues, while
rescaling only its largest eigenvalue to each case, predicts 16,013 updates
at every gamma under the saved step rule. The observed separation requires
deformation of the target-weighted normalized spectrum. All four runs attain
1\%, so failure to represent that tolerance does not explain their different times.

The other archived targets are a single sine, an exponential of sine, a
Runge function, and a quadratic. Across all five targets and four gammas,
all 19 observed hits lie in their computed intervals; the quadratic/gamma-12
run stopped at 200,000 updates and is censored. Only the primary sine-mixture
endpoints have the independent interval audit described next. These are
retrospective training-optimization comparisons, not held-out generalization results.

\section{What exactly was certified numerically?}
The analytic proofs above are exact-arithmetic statements. FP64 feature
assembly, SVD reconstruction, and action estimates require additional numerical
error control. The archived calculation records sensitivity allowances, but
these are not automatically interval enclosures of every plotted quantity.
In particular, Figure 1's masses and cutoff bounds have not received an
independent interval certificate.

The primary timing endpoints were separately verified using 192-bit Arb ball
arithmetic on nominal real tanh, with the archived binary inputs, centers,
target values, and step interpreted exactly. The independent checker uses
neither the SVD nor executed GD trajectories. Its algebra is as follows.

Put $H=J_\gamma^TJ_\gamma$ and $g=J_\gamma^Ty$. The coefficient recurrence is
$\theta_{n+1}=(I-\eta_\gamma H)\theta_n+\eta_\gamma g$. Hence, defining
\begin{equation*}
 \mathcal B=\begin{pmatrix}I-\eta_\gamma H&\eta_\gamma g\\0&1\end{pmatrix},
 \qquad \begin{pmatrix}\theta_n\\1\end{pmatrix}
 =\mathcal B^n\begin{pmatrix}0\\1\end{pmatrix},
\end{equation*}
gives the exact identity
\begin{equation}\label{eq:audit}
 E_\gamma(n)^2=
 \frac{\|y\|^2-2\theta_n^Tg+\theta_n^TH\theta_n}{\|y\|^2}.
\end{equation}
Both follow by induction and expansion of $\|y-J_\gamma\theta_n\|^2$.
Integer matrix powers evaluate \eqref{eq:audit} without stepping through an
optimizer trajectory.

To construct the common-slope Gram matrix efficiently, let $m_j$ be the
empirical mean of $\tanh(\gamma(x-c_j))$. For distinct centers,
\begin{equation*}
 H_{jk}=1-\frac{m_j-m_k}{\tanh(\gamma(c_k-c_j))}.
\end{equation*}
Indeed, the tanh subtraction identity gives
$a-b=\tanh(\gamma(c_k-c_j))(1-ab)$ for the two features at any input;
average and rearrange. Diagonal entries and repeated-center pairs use direct
products; $H_{00}=1$ and $H_{0j}=m_j$.

The audit first encloses $\eta_\gamma\|H\|$ below one using a row-sum bound.
For a proposed interval $[N_-,N_+]$, it then encloses \eqref{eq:audit} at the
excluded step $N_--1$ and sufficient step $N_+$. An enclosure strictly above the exact threshold $1/10000$
excludes the first, while an enclosure at or below it confirms the second.
Proposition 1's monotonicity completes the certificate. The reported
certificates concern these endpoints, not every FP64 curve, target, or newly
stated coarse cap bound. They also distinguish the nominal-real problem from
floating-point executed GD, whose measured crossings are compared separately.

\section{Bandwidth scaling and the final scope}
When centers have spacing $h$, keep $\lambda=\gamma h$ for dimensionless
bandwidth. At a frequency fixed relative to the grid, $\vartheta=\omega h$,
\begin{equation*}
 M_\gamma(\vartheta/h)=
 \frac{\pi|\vartheta|/(2\lambda)}{\sinh(\pi|\vartheta|/(2\lambda))}.
\end{equation*}
Substitution into \eqref{eq:multiplier} proves that fixed $\lambda$ preserves
attenuation at a fixed grid-relative frequency. For $h=2/N$ and
$W=N+2\lceil\sqrt N\rceil+1$, fixed positive $\lambda$ implies
$\gamma=\lambda N/2=\Theta(W)$. This does not imply constant training time
across widths: the target's sampled spectrum, geometry, and normalized rates
still enter \eqref{eq:gd}.

The complete logical chain is now explicit. Gamma controls the feature
multiplier. Theorem 1 bounds kernel action on a specified subspace and converts
target overlap with it into slow spectral mass and a delay. Theorem 2 retains
all geometric couplings and target weights, allowing much tighter,
case-specific time intervals. The empirical spectral plot checks the
intermediate allocation of target energy; the executed curves check its
learning consequences. No theorem here states that larger gamma helps every
target, that a gamma cap alone guarantees difficulty without a target
condition, or that the coarse subspace bound explains the measured factor
with the accuracy of the full-kernel calculation.

\begingroup\small
\section*{Sources and reproduction}
The proofs are self-contained apart from standard finite-dimensional spectral
theory, Fourier series, and the residue theorem. The classical GD spectral
interpretation is consistent with Y.~Yao, L.~Rosasco, and A.~Caponnetto,
\href{https://yao-lab.github.io/publications/YaoCapRos07_EarlyStop.pdf}{On Early
Stopping in Gradient Descent Learning}, \textit{Constructive Approximation}
26:289--315 (2007), \S3.3. The contribution being evaluated is the explicit
gamma-to-finite-kernel connection and its target-dependent quantitative consequences.

The \href{\evidence/REPORT.md}{study report},
\href{\evidence/interval_audit.json}{interval audit},
\href{\evidence/spectral_bridge_data.json}{spectral figure data}, and
\href{\evidence/pi_brief_figure_data.json}{executed-curve provenance} retain
the inputs and evidence roles. Earlier detailed constructions remain in
\href{\repo/docs/gamma_factorized_readout.md}{the technical note}; its
historical eigenvalue notation differs from this note.

Reproduce Figure 1 with the
\href{\repo/experiments/expD36_frozen_gamma_probe/spectral_bridge_figure.py}{spectral figure script}
and Figure 2 with the
\href{\repo/experiments/expD36_frozen_gamma_probe/pi_brief_figure.py}{GD/Adam figure script}.
Compile the directly authored
\href{\repo/docs/gamma_optimization_full_note.tex}{LaTeX source}
with \texttt{latexmk -pdf} from the repository root. The
\href{\repo/docs/gamma_optimization_full_note.md}{Markdown companion}
contains the full argument and commands. This write-up and spectral
diagnostic use the existing archive and no additional GPU training.
\endgroup
\end{document}
